\section{Topological Sort using Kahn's Algorithm (BFS)}
\label{sec:graph-topo-sort-kahn}

\problemstatement{Same as Section~\ref{sec:graph-topo-sort-dfs}: return
a valid topological ordering of a DAG -- this time using an
in-degree-based BFS instead of DFS finish times.}

\begin{patternbox}
\emph{``Repeatedly do whatever has no remaining prerequisites''} is
solved directly: track each node's \textbf{in-degree} (number of
unprocessed incoming edges); start with every in-degree-$0$ node in a
queue; each time a node is processed, decrement its neighbors'
in-degrees, and any neighbor that drops to $0$ joins the queue.
\end{patternbox}

\subsection*{Why this directly produces a valid order}

A node can only be safely placed in the output once \textbf{every}
prerequisite pointing into it has already been placed -- which is
exactly what in-degree $0$ means at the moment a node is dequeued (all
of its incoming edges have already been ``consumed'' by their source
nodes being processed earlier). This makes Kahn's algorithm a direct
simulation of the real-world process the topological sort is modeling:
do everything that's ready, which unlocks more things to become ready,
repeat.

\subsection*{Algorithm}

\begin{enumerate}
  \item Compute the in-degree of every node by scanning all edges.
  \item Push every node with in-degree $0$ into a queue.
  \item While the queue is non-empty: pop a node, append it to the
        result; for each of its neighbors, decrement their in-degree,
        and push any neighbor whose in-degree reaches $0$.
\end{enumerate}

\subsection*{C++ implementation}

\begin{lstlisting}
#include <bits/stdc++.h>
using namespace std;

vector<int> topoSortKahn(int n, vector<vector<int>>& adj) {
    vector<int> inDegree(n, 0);
    for (int u = 0; u < n; u++) {
        for (int v : adj[u]) inDegree[v]++;
    }

    queue<int> q;
    for (int i = 0; i < n; i++) {
        if (inDegree[i] == 0) q.push(i);
    }

    vector<int> result;
    while (!q.empty()) {
        int node = q.front();
        q.pop();
        result.push_back(node);

        for (int neighbor : adj[node]) {
            inDegree[neighbor]--;
            if (inDegree[neighbor] == 0) q.push(neighbor);
        }
    }
    return result;
}

int main() {
    int n = 6;
    vector<vector<int>> adj(n);
    adj[5] = {2, 0}; adj[4] = {0, 1}; adj[2] = {3}; adj[3] = {1};

    for (int x : topoSortKahn(n, adj)) cout << x << " ";
    cout << endl; // 4 5 2 0 3 1  (one valid order)
    return 0;
}
\end{lstlisting}

\subsection*{Dry run}

\dryrun{Same graph as Section~\ref{sec:graph-topo-sort-dfs}: edges
$5{\to}2,5{\to}0,4{\to}0,4{\to}1,2{\to}3,3{\to}1$.}

In-degrees: $0$ has in-degree $2$ (from $5,4$); $1$ has in-degree $2$
(from $4,3$); $2$ has in-degree $1$ (from $5$); $3$ has in-degree $1$
(from $2$); $4,5$ have in-degree $0$. Initial queue: $[4,5]$. Pop $4$:
result $[4]$; decrement $0\to1$, $1\to1$ (neither reaches $0$ yet).
Pop $5$: result $[4,5]$; decrement $2\to0$ (push $2$), $0\to0$ (push
$0$). Queue: $[2,0]$. Pop $2$: result $[4,5,2]$; decrement $3\to0$
(push $3$). Pop $0$: result $[4,5,2,0]$; no outgoing edges. Pop $3$:
result $[4,5,2,0,3]$; decrement $1\to0$ (push $1$). Pop $1$: result
$[4,5,2,0,3,1]$. Final order: $4,5,2,0,3,1$ -- a different but equally
valid ordering from the DFS approach's $5,4,2,3,1,0$ (topological
order is rarely unique).

\begin{complexitybox}
\textbf{Time Complexity.} $O(V+E)$ -- computing in-degrees scans every
edge once; the BFS itself also visits every node and edge once.
\textbf{Space Complexity.} $O(V)$ for the in-degree array and the
queue.
\end{complexitybox}

\begin{takeawaybox}
Kahn's algorithm has one advantage DFS's approach doesn't offer for
free: if the queue empties before every node has been processed, the
remaining unprocessed nodes are necessarily stuck in a cycle (nothing
in a cycle can ever reach in-degree $0$, since each node in the cycle
is always waiting on another node in the same cycle). This makes Kahn's
algorithm double as a \textbf{cycle detector} in directed graphs, used
directly in the next section.
\end{takeawaybox}
